\chapter{Implementation}

\section{Technology Stack}
\begin{table}[H]
\centering
\caption{Technologies used.}
\small
\begin{tabularx}{\textwidth}{lX}
\toprule
Area & Tools \\
\midrule
Language & Python 3.11 \\
Data & Objaverse API, pandas \\
3D processing & trimesh~\cite{trimesh}, NumPy, SciPy (KD-tree, flood fill), scikit-image (marching cubes) \\
Rendering & Blender 4.2 LTS~\cite{blender}, Cycles on GPU \\
Learning & PyTorch~\cite{paszke2019pytorch}, timm (DINOv2) \\
Simplification & own QEM implementation; PyMeshLab as baseline only \\
Demo & Gradio, rembg \\
Compute & Kaggle notebooks (4-core CPU, 2$\times$ NVIDIA T4 GPU), Kaggle API \\
Version control & Git, GitHub \\
\bottomrule
\end{tabularx}
\end{table}

\section{Code Structure}
\begin{table}[H]
\centering
\caption{Main modules.}
\small
\begin{tabularx}{\textwidth}{lX}
\toprule
Path & Responsibility \\
\midrule
\texttt{configs/default.yaml} & every setting (filters, grid size, model, training); no hard-coded values \\
\texttt{gen3d/data/mesh\_io.py} & load GLB, normalise, sample coloured surface points \\
\texttt{gen3d/data/occupancy.py} & occupancy labelling (Algorithm~\ref{alg:occ}) \\
\texttt{gen3d/data/dataset.py} & PyTorch dataset: image, camera, points; $1/\sqrt{n}$ sampler \\
\texttt{gen3d/generator/encoder.py} & DINOv2 encoder (ResNet-18 kept for ablation) \\
\texttt{gen3d/generator/model.py} & generator, loss, IoU \\
\texttt{gen3d/generator/extract.py} & mesh extraction \\
\texttt{scripts/*.py} & pipeline steps: select, download, check, preprocess, render, review, train, generate \\
\texttt{kaggle/} & builds and pushes the Kaggle notebooks (data, render, train) \\
\bottomrule
\end{tabularx}
\end{table}

\section{Cloud Execution on Kaggle}
The laptop has no GPU, so all heavy work runs on Kaggle and is controlled from the laptop through the Kaggle API: the code is uploaded as a private dataset, each step is pushed as a notebook, and outputs are downloaded for inspection. Each notebook processes objects in chunks, saves every finished chunk immediately, stops before Kaggle's 12-hour limit, and resumes from where it stopped. The data notebook downloads, checks and preprocesses objects; the render notebook installs Blender, re-downloads each chunk of models, renders them with several Blender processes on two GPUs and stores each object's views as one zip file; the training notebook reads the outputs of both.

\section{Coding Conventions}
\begin{itemize}[leftmargin=*]
  \item Modular package (\texttt{gen3d/}) separated from runnable steps (\texttt{scripts/}).
  \item All parameters in one YAML configuration file.
  \item Docstrings on every module and function; descriptive names.
  \item Core algorithms (occupancy labelling, QEM, metrics) implemented by the team rather than taken from libraries.
  \item Every step is resumable and logs its decisions to CSV files, which also serve as the data dictionary.
\end{itemize}
